% Capítulo: Evaluation
% Codificación: UTF-8

\chapter{Evaluation}
\label{chap:evaluation}

The specification is evaluated with a test suite. The suite has two
parts, each with its own reference. The smoke tests are handwritten
kernels, one per instruction, that exercise all of its variants and
compare the results with golden values derived from the PISA
document~\cite{pisa-spec}. The kernel tests run whole kernels compiled
from OpenCL and compare their outputs with an expected output computed
from the OpenCL source, by a sequential C++ program or by hand.
Every test produces a trace of the execution, step by step, which is how
the specification was debugged and where the traces shown in this
chapter come from. The suite ran after every new feature and caught the
bugs of the first versions, mainly in the handling of divergence.

As can be expected, the suite does not cover every possible case of the
subset. It is complete enough to show that execution within a sub-group
is correct: that each instruction computes what PISA prescribes, and
that kernels compose them correctly, including kernels whose work-items
diverge.

\section{Smoke Tests}
\label{sec:eval-smoke}

Since the specification was designed so that the semantics of each
instruction is a function, the smoke tests check that each of these
functions computes what the PISA document describes. There is one smoke
test per instruction, a handwritten kernel that covers every variant of
the instruction, the edge cases of its definition and the examples the
document gives.

A smoke test is a kernel with no arguments, followed by the golden values
it must produce. For each case, the kernel sets up the operands with a
few instructions, executes the subject instruction variant, and moves on to
the next case. The golden values are the expected
contents of each destination for each lane, derived from the definition
in the document. The kernel goes through the parser and runs in ASLi,
with all 32 lanes active and no host library involved. A script then
reads the values written from the trace and compares them with the
golden ones. Listing~\ref{lst:smoke-iadd} shows one case from the smoke
test of \texttt{iadd}: each lane adds $-16$ to its lane id, so lanes 0
to 15 obtain a negative result and lanes 16 to 31 a non-negative one, which
checks that the addition is signed and wraps. The suite has 26 smoke tests, 
one per instruction, with around 300 cases in all, and every case matches 
its golden values.

\begin{lstlisting}[caption={[Smoke test for iadd]Smoke
test for \texttt{iadd.32b}. The kernel is followed by the golden values
of its destination register, one per lane.}, label={lst:smoke-iadd}]
.kernel @iadd()
{
  .reg .32b %r0
  iadd.32b %r0, %laneid, -16
  return
}
GOLD
%r0 : 0xfffffff0 0xfffffff1 0xfffffff2 ... 0xd 0xe 0xf
\end{lstlisting}

When implementing \texttt{shfl}, the smoke tests made it possible to
identify two errors in the examples of the PISA
document~\cite{pisa-spec}. One was a prefix-sum example that assumed
lanes whose source is out of range do not participate, while the
actual definition says they retain their own value. The other tried to show a
segmented shuffle that moved its items down, but set its clamp value to
0, so once computed no segment was shuffled at all. Executing the
examples as written exposed both errors, and they were reported
to the PISA maintainers.

With this, the semantics of each instruction can be taken as correct.
What remains is to see how they behave together in real workloads, under
divergence, and with data that comes from a host program.


\section{Kernel Tests}
\label{sec:eval-results}

The kernel tests are built on the host library. Each test is a kernel
written in OpenCL C and compiled to PISA with a pre-release build of the
LLVM PISA compiler. Each test also has a host program that sets the
NDRange, fills the input buffers, launches the kernel through the
library, and reads the output buffer back. The expected
output of each kernel is given to the test, and the test passes when the
buffer read back matches it element by element.

The suite has 23 kernels, listed in Table~\ref{tab:suite} with their
sources in Appendix~\ref{app:kernels}, in four groups. The first group checks the basics: a kernel with no branches,
launched over one and over several work-groups, to see that arguments,
buffers and the global id reach the kernel and the results come back.
The control flow group checks that the divergence model holds from the
simple cases to the complex ones. It starts with one branch where the 
lanes split and rejoin, goes on to
nested branches and to divergent loops, and ends with two kernels that
mix loops, switches, breaks and continues into the most complex control
flow of the suite. The sub-group
communication group runs the four shuffle patterns and the sum
reduction, with all lanes active and inside divergent regions, over
one and over several sub-groups. The memory group checks that the lanes
of a work-group reach one shared
space, where a value stored by one lane can be loaded by another: each
lane stores its element and loads another lane's in a later step, with
no barrier between the two, since the subset has none. In OpenCL that is
a race. Under lockstep execution the store of one step is visible to the
load of the next, so the kernel's output is defined by the execution
model, and the exchange stays within a sub-group.

\begin{table}[t]
\centering
\small
\begin{tabular}{lp{0.55\linewidth}>{\raggedright\arraybackslash}p{0.15\linewidth}}
\toprule
Kernel & What it exercises &  Work-groups $\times$ sub-groups \\
\midrule
\multicolumn{3}{l}{\emph{Basics}} \\
\texttt{vadd}            & straight-line code, global loads and stores & 1$\times$1, 2$\times$2 \\
\midrule
\multicolumn{3}{l}{\emph{Control flow}} \\
\texttt{vadd\_if\_else}  & one divergent branch, both paths store & 1$\times$1 \\
\texttt{vadd\_mod}       & one divergent branch, half the lanes skip a region and store nothing & 1$\times$1 \\
\texttt{nested\_if}      & tree of nested divergent branches, one join before the store & 1$\times$1 \\
\texttt{loop}            & loop with lane-dependent trip count & 1$\times$2 \\
\texttt{nested\_loop}    & nested loop inside a lane-dependent loop & 1$\times$2 \\
\texttt{control\_flow}   & lane-dependent loop around a switch with fall-through, break, continue and a do-while & 1$\times$1 \\
\texttt{control\_flow1}  & as \texttt{control\_flow}, with a nested switch and an inner loop & 1$\times$1 \\
\midrule
\multicolumn{3}{l}{\emph{Sub-group communication}} \\
\texttt{shfl1}           & \texttt{shfl.idx} inside nested uniform loops, sum over lanes & 1$\times$1 \\
\texttt{shfl2}           & \texttt{shfl.idx} inside a uniform loop, broadcast of lane 31 & 1$\times$1 \\
\texttt{shfl3}           & \texttt{shfl.idx} from the next lane, index computed with \texttt{sdiv.rem} & 1$\times$1 \\
\texttt{shfl\_dn1}       & rotate-down reduction & 1$\times$1 \\
\texttt{shfl\_dn2}       & zero-fill down-reduction, unrolled, inside a uniform loop, only lane 0 stores & 1$\times$1 \\
\texttt{shfl\_dn3}       & as \texttt{shfl\_dn2}, with the reduction as a loop whose bound is computed at run time & 1$\times$1 \\
\texttt{shfl\_up1}       & rotate-up reduction, unrolled & 1$\times$1 \\
\texttt{shfl\_up2}       & rotate-up reduction, as a loop & 1$\times$1 \\
\texttt{shfl\_up3}       & zero-fill up-reduction, unrolled, inside a uniform loop, only lane 31 stores & 1$\times$1 \\
\texttt{shfl\_xor1}      & lane reversal with \texttt{shfl.xor} & 1$\times$1 \\
\texttt{shfl\_xor2}      & reversal with \texttt{shfl.xor} in a divergent region, lanes 0 to 15 only & 1$\times$1 \\
\texttt{prefix\_sum}     & inclusive scan, compiled to shuffles up & 1$\times$1, 2$\times$1 \\
\texttt{ired\_sum}       & \texttt{ired.sum} over the sub-group & 1$\times$1, 1$\times$2 \\
\texttt{reduce\_if\_else}& \texttt{ired.sum} in a divergent region, lanes 0 to 15 only & 1$\times$1, 1$\times$2 \\
\midrule
\multicolumn{3}{l}{\emph{Memory}} \\
\texttt{exchange}        & values passed between the lanes of a sub-group through shared memory, with no barrier & 2$\times$2 \\
\bottomrule
\end{tabular}
\caption[The kernel suite]{The kernel suite.}
\label{tab:suite}
\end{table}

Every kernel is launched over work-groups whose size is a multiple of
32, and the last column of Table~\ref{tab:suite} shows which kernels
also run over several work-groups or several sub-groups per work-group.
The specification records a trace of each sub-group: for every step,
the global program counter, the active mask, the instruction in PISA
syntax, and the values written. The traces in this section are
abridged to the steps that matter, and per-lane values are shown as the
pattern they follow.

Every kernel matched its expected output at every work-item, in every
launch. For a kernel without divergence, a matching output is enough.
For a divergent kernel it is not: a wrong schedule can still produce
the right output when the paths do not interact. The traces are needed
to see that the lanes diverge and reconverge as the model specifies.
The rest of this section follows four kernels through their traces: a
kernel without divergence, a kernel that diverges forward through
nested branches, a loop whose lanes leave at different times, and a
reduction under a partial mask. Masks are written with lane~0 as the
least significant bit, and a step at a label is shown by the label's
name.

\paragraph{Uniform Execution: \texttt{vadd}.}
The first kernel adds two vectors, \texttt{c[i] = a[i] + b[i]}, and has
no divergence: every work-item executes the same instructions, and the
mask is \texttt{ffffffff} at every step. It shows that the kernel as a
whole computes what its OpenCL source describes, and that data
submitted by the host reaches the kernel and comes back.

\begin{lstlisting}[language=C, caption={[Source of vadd]OpenCL source
of \texttt{vadd}.}, label={lst:src-vadd}]
__kernel void vadd(__global int *a, __global int *b, __global int *c)
{
    int i = get_global_id(0);
    c[i] = a[i] + b[i];
}
\end{lstlisting}

\begin{lstlisting}[caption={[Trace of vadd]Trace of \texttt{vadd},
abridged.}, label={lst:trace-vadd}]
PC: 0x0   mask: ffffffff   ld.param.64b %r0, [%arg0]
          %r0: 0x0 in every lane
PC: 0x1   mask: ffffffff   ld.param.64b %r1, [%arg1]
          %r1: 0x80 in every lane
PC: 0x2   mask: ffffffff   ld.param.64b %r2, [%arg2]
          %r2: 0x100 in every lane
PC: 0x3   mask: ffffffff   zext.32b.64b %r3, %localid.x
          %r3: 0, 1, 2, ..., 31
...
PC: 0x7   mask: ffffffff   asr.64b %r7, %r6, 30
          %r7: 0x0, 0x4, 0x8, ..., 0x7c
PC: 0x8   mask: ffffffff   ld.global.32b %w0, [%r0 + %r7]
          %w0: 1, 2, 3, ..., 32
PC: 0x9   mask: ffffffff   ld.global.32b %w1, [%r1 + %r7]
          %w1: 1, 2, 3, ..., 32
PC: 0xa   mask: ffffffff   iadd.32b %w2, %w1, %w0
          %w2: 2, 4, 6, ..., 64
PC: 0xb   mask: ffffffff   st.global.32b [%r2 + %r7], %w2
PC: 0xc   mask: ffffffff   return
\end{lstlisting}

Three things in the trace are worth noting. The three buffer addresses
read from parameter memory, \texttt{0x0}, \texttt{0x80} and
\texttt{0x100}, are 128 bytes apart, one buffer of 32 four-byte
elements after another, as the host library laid them out. The byte
offset in \texttt{\%r7} is four times the work-item's id, so each lane
addresses its own element. And the values stored at \texttt{0xb} run
2, 4, 6, up to 64, which is what the host reads back.

\paragraph{Forward Divergence: \texttt{nested\_if}.}
The second kernel loads one value per work-item and runs it through
nested conditionals, writing a result that depends on the path taken.
In the trace the value is the work-item's id, so the branches split the
sub-group into contiguous halves and quarters. The two innermost
levels compile to \texttt{sel}, so the kernel diverges three
times, into four groups of eight lanes, each with its own block, and
all four rejoin at one label before the store. This is the
\texttt{if}-\texttt{else} of Table~\ref{tab:divergence-if-else},
nested.

\begin{lstlisting}[caption={[Trace of nested\_if]Trace of
\texttt{nested\_if}, abridged: the three splits, the four blocks, and
the join.}, label={lst:trace-forward}]
PC: 0x5    mask: ffffffff   scmp.gt.32b %p0, %r4, 15
PC: 0x6    mask: ffffffff   goto.cond %p0, LBB0_4      ; lanes 16-31 leave
PC: 0x7    mask: 0000ffff   scmp.gt.32b %p1, %r4, 7
PC: 0x8    mask: 0000ffff   goto.cond %p1, LBB0_3      ; lanes 8-15 leave
PC: 0x9    mask: 000000ff   scmp.gt.32b %p2, %r4, 3    ; block for lanes 0-7
...
PC: 0x10   mask: 000000ff   goto LBB0_7
PC: 0x11   mask: 0000ff00   LBB0_3:                     ; lanes 8-15
...
PC: 0x18   mask: 0000ff00   goto LBB0_7
PC: 0x19   mask: ffff0000   LBB0_4:                     ; lanes 16-31
PC: 0x1a   mask: ffff0000   ucmp.gt.32b %p8, %r4, 23
PC: 0x1b   mask: ffff0000   goto.cond %p8, LBB0_6      ; lanes 24-31 leave
PC: 0x1c   mask: 00ff0000   ucmp.gt.32b %p9, %r4, 19   ; block for lanes 16-23
...
PC: 0x22   mask: 00ff0000   goto LBB0_7
PC: 0x23   mask: ff000000   LBB0_6:                     ; lanes 24-31
...
PC: 0x29   mask: ff000000   sel.16b %r8, %r13, %r14, %p12
PC: 0x2a   mask: ffffffff   LBB0_7:                     ; all rejoin
PC: 0x2b   mask: ffffffff   zext.16b.32b %r15, %r8
           %r15: 0, 0, 2, 2, 4, 4, ..., 30, 30
PC: 0x2c   mask: ffffffff   st.global.32b [%r1 + %r3], %r15
\end{lstlisting}

The point of the trace is the order of the blocks and the masks they
run under. After the splits, lanes 0 to 7 run their block alone while
lanes 8 to 15 wait at \texttt{0x11}, lanes 16 to 31 at \texttt{0x19},
and the finished lanes at the join. The lowest counter runs next, so
the four blocks execute in address order, each under the mask of its
own eight lanes, and no lane runs a block that is not its own. When
the last block reaches \texttt{0x2a}, every counter is at the join,
the mask returns to \texttt{ffffffff}, and the store executes for the
whole sub-group with each lane holding the value its own path
computed.

\paragraph{Backward Divergence: \texttt{loop}.}
The third kernel is a loop that each lane runs a different number of
times: lane $i$ runs it $i$ times and accumulates the first $i$
elements of its sub-group's block (Listing~\ref{lst:src-loop}).
Lane 0 has nothing to do and goes straight to the exit at
\texttt{LBB0\_3}, so the mask is \texttt{fffffffe} from the first
instruction after the branch. The body ends with the back-edge at
\texttt{0x13}, taken while the lane's counter is below its lane id. At
the end of every trip one lane has reached its count, falls through to
the exit, and waits there while the others go round again. The lanes
leave one per trip, in order of their id, and
Table~\ref{tab:mask-staircase} shows the mask at the back-edge on
successive trips: the cleared bits grow from the right, one per trip,
until only lane 31 remains. The only branch inside the loop is the
back-edge, so the body itself never diverges. When lane 31 falls
through, all 32 lanes are at the exit, the mask returns to
\texttt{ffffffff}, and the store runs once for every lane. The whole
kernel takes 296 steps, 31 of them back-edges. The second sub-group of
the launch produces a trace of the same shape, over its own block of
the input.

\begin{lstlisting}[language=C, caption={[Source of loop]OpenCL source
of \texttt{loop}.}, label={lst:src-loop}]
__kernel void loop(__global const int *a, __global int *out)
{
    int gid  = get_global_id(0);
    int lid  = get_sub_group_local_id();
    int base = gid - lid;
    int s = 0;
    for (int i = 0; i < lid; i++)
        s += a[base + i];
    out[gid] = s;
}
\end{lstlisting}

\begin{lstlisting}[caption={[Trace of loop]Trace of \texttt{loop},
abridged: the first back-edge, the last one, and the exit.},
label={lst:trace-loop}]
PC: 0x6   mask: ffffffff   ucmp.eq.32b %p0, %laneid, 0
PC: 0x7   mask: ffffffff   goto.cond %p0, LBB0_3
PC: 0x8   mask: fffffffe   trunc.64b.32b %w1, %d4
PC: 0x9   mask: fffffffe   isub.32b %w2, %w1, %laneid
PC: 0xa   mask: fffffffe   mov.32b %w3, %w0
PC: 0xb   mask: fffffffe   LBB0_2:
PC: 0xc   mask: fffffffe   iadd.32b %w4, %w2, %w3
PC: 0xd   mask: fffffffe   sext.32b.64b %d5, %w4
PC: 0xe   mask: fffffffe   shl.64b %d6, %d5, 2
PC: 0xf   mask: fffffffe   ld.global.32b %w5, [%d0 + %d6]
PC: 0x10  mask: fffffffe   iadd.32b %w0, %w5, %w0
PC: 0x11  mask: fffffffe   iadd.32b %w3, %w3, 1
PC: 0x12  mask: fffffffe   ucmp.lt.32b %p1, %w3, %laneid
PC: 0x13  mask: fffffffe   goto.cond %p1, LBB0_2
PC: 0xb   mask: fffffffc   LBB0_2:
...
PC: 0x13  mask: fffffffc   goto.cond %p1, LBB0_2
PC: 0xb   mask: fffffff8   LBB0_2:
...
PC: 0x13  mask: 80000000   goto.cond %p1, LBB0_2
PC: 0x14  mask: ffffffff   LBB0_3:
PC: 0x15  mask: ffffffff   shl.64b %d7, %d4, 32
PC: 0x16  mask: ffffffff   asr.64b %d8, %d7, 32
PC: 0x17  mask: ffffffff   shl.64b %d9, %d8, 2
PC: 0x18  mask: ffffffff   st.global.32b [%d1 + %d9], %w0
PC: 0x19  mask: ffffffff   return
\end{lstlisting}

In Listing~\ref{lst:trace-loop}, each omitted block is one trip of the body, from the
label at \texttt{0xb} to the back-edge, under a mask with one more bit
cleared. After the last back-edge the active lanes are exactly the
lanes at the exit label, all 32 of them.

\begin{table}[t]
\centering
\begin{tabular}{ccc}
\toprule
Trip & Mask at back-edge & Active lanes \\
\midrule
1  & \texttt{fffffffe} & 31 \\
2  & \texttt{fffffffc} & 30 \\
3  & \texttt{fffffff8} & 29 \\
\multicolumn{3}{c}{$\vdots$} \\
29 & \texttt{e0000000} & 3 \\
30 & \texttt{c0000000} & 2 \\
31 & \texttt{80000000} & 1 \\
\bottomrule
\end{tabular}
\caption[Active mask at the loop back-edge]{Active mask at the back-edge
of \texttt{loop} on successive trips.}
\label{tab:mask-staircase}
\end{table}

\paragraph{Communication under Divergence: \texttt{reduce\_if\_else}.}
The fourth kernel sums the input over the sub-group with
\texttt{ired.sum}, but only in lanes 0 to 15, which take one path of a
branch. The other 16 lanes take the other path and write $-1$
(Listing~\ref{lst:src-reduce}). The result of the reduction depends
on which lanes are active when it executes, so this is one of the
kernels on which the reconvergence model can be observed through the
output.

\begin{lstlisting}[language=C, caption={[Source of reduce\_if\_else]OpenCL
source of \texttt{reduce\_if\_else}.}, label={lst:src-reduce}]
__kernel void reduce_if_else(__global const int *a, __global int *out)
{
    int gid = get_global_id(0);
    int lid = get_sub_group_local_id();
    if (lid < 16)
        out[gid] = sub_group_reduce_add(a[gid]);
    else
        out[gid] = -1;
}
\end{lstlisting}

\begin{lstlisting}[caption={[Trace of reduce\_if\_else]Trace of
\texttt{reduce\_if\_else}, abridged: the branch, the two paths, and
the join.}, label={lst:trace-reduce}]
PC: 0x5   mask: ffffffff   ucmp.lt.32b %p0, %laneid, 16
PC: 0x8   mask: ffffffff   goto.cond %p0, LBB0_2      ; lanes 0-15 leave
PC: 0x9   mask: ffff0000   mov.32b %w0, -1            ; lanes 16-31
PC: 0xa   mask: ffff0000   goto LBB0_3
PC: 0xb   mask: 0000ffff   LBB0_2:                    ; lanes 0-15
PC: 0xc   mask: 0000ffff   ld.global.32b %w1, [%d0 + %d6]
          %w1: 1, 2, 3, ..., 16 in lanes 0-15
PC: 0xd   mask: 0000ffff   ired.sum.32b %w0, %w1, -1
          %w0: 136 in lanes 0-15, -1 in lanes 16-31
PC: 0xe   mask: ffffffff   LBB0_3:                    ; all rejoin
PC: 0xf   mask: ffffffff   st.global.32b [%d1 + %d6], %w0
PC: 0x10  mask: ffffffff   return
\end{lstlisting}

Two things in the trace (Listing~\ref{lst:trace-reduce}) matter. The
third operand of \texttt{ired.sum} is $-1$, every lane, and yet the
result is 136, the sum of 1 to 16: PISA defines the reduction over the
participating lanes that are active, and the model decides which lanes
are active at \texttt{0xd}: the 16 that took the branch. And the two paths
run in address order, lanes 16 to 31 first at \texttt{0x9} and
\texttt{0xa}, lanes 0 to 15 next at \texttt{0xb} to \texttt{0xd}, with
the join at \texttt{0xe} once every counter is there. The buffer the
host reads back holds 136 in its first 16 elements and $-1$ in
the rest, which is the expected output.

In the three divergent kernels the trace shows what the worked examples
of Section~\ref{sec:asl-divergence} showed on six instructions: the
active mask and the global program counter follow from the per-lane
counters alone, the lanes furthest behind run first, and reconvergence
happens at the step where the counters coincide, with no reconvergence
point recorded anywhere.

For 21 of the 23 kernels, the output is fixed by the instruction
semantics alone, so any correct execution of the subset must produce
it. The other two, \texttt{shfl\_xor2} and \texttt{reduce\_if\_else},
run a shuffle or a reduction under a partial mask, and their output
depends on which lanes are active at that instruction, which PISA does
not define. For these two, the expected output assumes that the lanes on the same
path reach the instruction together, as they would in a lockstep
execution of structured code. The specification produces that output,
so its reconvergence rule
activates the expected lanes.

Taken together, the smoke tests show that each instruction of the
subset computes what PISA prescribes, the kernel tests show that the
instructions compose and that the sub-group diverges and reconverges as
the model specifies.
